\documentclass[a4paper,12pt]{article}

\title{Blocco 3}
\author{Leonardo Centazzo}
\date{\today}


% Pacchetti essenziali
\usepackage{amsmath, amssymb, amsthm}
\usepackage{enumitem}
\usepackage{geometry}

% Impostazioni pagina minime
\geometry{a4paper, margin=2cm}

% Stili per teoremi, proposizioni, ecc.
\theoremstyle{plain}
\newtheorem{teorema}{Teorema}
\newtheorem{proposizione}{Proposizione}
\newtheorem{lemma}{Lemma}
\newtheorem{corollario}{Corollario}

\theoremstyle{definition}
\newtheorem{definizione}[teorema]{Definizione}
\newtheorem{esempio}[teorema]{Esempio}
\newtheorem{esercizio}{Esercizio}
\newtheorem{osservazione}[teorema]{Osservazione}

\newtheorem*{warning}{Warning}

\theoremstyle{remark}
\newtheorem{nota}[teorema]{Nota}


\begin{document}

\maketitle

In the following hf will stand for hyperfinite and hhf will stand for hyperhyperfinite.

\begin{esercizio}

	Define an equivalence relation on the Baire space by setting
	$x E y$ iff for all large $n$, $x_n = y_n$. Prove that $E$ is hf.

\end{esercizio}

\begin{proof}
	

\end{proof}





\begin{esercizio}

	Show that $E \leq_B E'$ and $E'$ is hhf, then $E$ is also hhf.

\end{esercizio}

\begin{proof}

	Since $E'$ is hf, we can write $E' = \bigcup_{n < \omega} E'_n$ with $<\langle E_n : n < \omega \rangle$ increasing sequence of hf cbers.
	Let $f: X \to X$ be a Borel reduction from $E$ to $E'$. 
	For every $n < \omega$ define $x E_n y \iff f(x) E'_n f(y)$.
	\begin{itemize}

		\item $E \subseteq \bigcup_{n < \omega} E_n$ :

			\begin{align}
				x E y \Rightarrow f(x) E' f(y) \Rightarrow f(x) E'_n f(y) \text{ for some }n < \omega
				\Rightarrow x E_n y \Rightarrow x \bigcup_{n < \omega} E_n y
			\end{align}

		\item $\bigcup_{n < \omega} E_n \subseteq E$ :

			\begin{align}
				x \bigcup_{n < \omega} E_n \Rightarrow x E_n y \text{ for some }n < \omega \Rightarrow
				f(x) E'_n f(y) \text{ for some }n < \omega \Rightarrow f(x) E' f(y) \Rightarrow x E y
			\end{align}

		\item for all $n < \omega \ E_n$ is hf cber : 
			
			Fix $n < \omega$. Then $E_n$ is a cber because $f$,$E'_n$ are Borel and $E_n \subseteq E$.
			$E_n$ is hf because by definition $f$ is a Borel reduction from $E_n$ to $E'_n$ and $E'_n$ is Borel.

		\item $\langle E_n : n < \omega \rangle$ is increasing :

			In fact $\langle E'_n : n < \omega \rangle$ is increasing, so we get
			\begin{align}
				x E_n y \Rightarrow f(x) E'_n f(y) \Rightarrow f(x) E'_{n+1} f(y) \Rightarrow f(x) E' f(y) 
				\Rightarrow x E y
			\end{align}

	\end{itemize}

	We conclude observing that $E = \bigcup_{n < \omega} E_n$ witnesses $E$ hf.

\end{proof}




\begin{esercizio}

	Show that the Cohen forcing $Add(\omega,1)$ is weakly homogeneous.

\end{esercizio}

\begin{proof}
\end{proof}




For the next problem, let $k$ be a measurable cardinal with a normal measure $U$ and let $P$ be the Prickry forcing at $k$.

\begin{esercizio}

	Prove that $P$ is weakly homogeneous.

\end{esercizio}

\begin{proof}
	

\end{proof}


\end{document}

